%=======================================================================
%  41277 Control Design -- Spring 2026
%  Demonstration Practical 2: State-Feedback Control
%  LaTeX version of the official Word answer template.
%
%  Compile with tectonic (see .vscode/settings.json) or pdflatex (twice).
%  If the ebgaramond package is unavailable, comment out the
%  \usepackage[lining]{ebgaramond} line below -- everything else still works.
%=======================================================================
\documentclass[11pt,a4paper]{article}

\usepackage[T1]{fontenc}
\usepackage[utf8]{inputenc}
\usepackage[lining]{ebgaramond}          % Garamond, to match the Word template
\usepackage[margin=2.2cm,top=3.2cm,headheight=42pt]{geometry}
\usepackage{amsmath,amssymb,bm}
\usepackage{mathtools}
\usepackage{empheq}
\usepackage{graphicx}
\usepackage{float}          % provides the [H] "exactly here" placement
\usepackage{subcaption}     % side-by-side plots and zooms
\usepackage{array}
\usepackage{booktabs}
\usepackage{fancyhdr}
\usepackage{enumitem}
\usepackage{xcolor}
\usepackage{listings}       % MATLAB snippets
\usepackage[colorlinks=true,linkcolor=black,urlcolor=blue!60!black,
            citecolor=black]{hyperref}

%----------------------------- house style -----------------------------
\setlength{\parindent}{0pt}
\setlength{\parskip}{6pt}
\allowdisplaybreaks
\raggedbottom

\newcommand{\dd}[1]{\dot{#1}}
\newcommand{\ddd}[1]{\ddot{#1}}
\newcommand{\todo}[1]{\textcolor{red}{\textbf{[TODO:} #1\textbf{]}}}

\lstset{language=Matlab,basicstyle=\ttfamily\small,frame=single,
        keywordstyle=\color{blue!70!black},commentstyle=\color{green!45!black},
        stringstyle=\color{purple!70!black},breaklines=true,columns=fullflexible}

% Section headings styled like the "Q1:", "Q2.1)" labels of the template
\usepackage{titlesec}
\titleformat{\section}{\large\bfseries}{}{0pt}{}
\titleformat{\subsection}{\normalsize\bfseries}{}{0pt}{}
\titlespacing*{\section}{0pt}{14pt}{4pt}
\titlespacing*{\subsection}{0pt}{10pt}{2pt}

% ---- figure that degrades to a placeholder box if the file is missing ----
%  \optfig{path}{width fraction}{caption}{label}
\newcommand{\optfig}[4]{%
  \begin{figure}[H]\centering
  \IfFileExists{#1}{\includegraphics[width=#2\linewidth]{#1}}%
    {\fbox{\parbox[c][3.2cm][c]{0.9\linewidth}{\centering
      \textit{[ Insert \texttt{\detokenize{#1}} here ]}}}}%
  \caption{#3}\label{#4}\end{figure}}

%------------------------------- header --------------------------------
\fancypagestyle{uts}{%
  \fancyhf{}
  \fancyhead[L]{\includegraphics[height=22pt]{assets/logos/uts_logo.png}\quad
                \textbf{Demonstration Practical 2}}
  \fancyhead[R]{\footnotesize\begin{tabular}{@{}r@{}}
                41277 Control Design\\ Spring 2026\\
                Coordinator: A/Prof Ricardo P. Aguilera
                \end{tabular}}
  \renewcommand{\headrulewidth}{0.4pt}
  \fancyfoot[C]{\thepage}
}
\pagestyle{uts}

%=======================================================================
\begin{document}

%=========================== COVER PAGE ================================
\thispagestyle{empty}

\includegraphics[height=30pt]{assets/logos/uts_logo.png}\hfill
{\large\bfseries Faculty of Engineering and IT}

\vspace{10pt}

\fbox{%
\begin{minipage}{\dimexpr\textwidth-2\fboxsep-2\fboxrule\relax}
\vspace{4pt}
\centering{\Large\bfseries Assignment Cover Sheet}\par
\vspace{8pt}
\raggedright
\textbf{Subject:} 41277 Control Design\par
\vspace{4pt}
\textbf{Assignment Number:} Demonstration Practical 2
\hfill \textbf{Date Submitted:} \todo{dd/mm/2026}\par
\vspace{4pt}
\textbf{Assignment Title:} State-Feedback Control\par
\vspace{4pt}
\textbf{Student Name(s) and Number(s)}\hfill\textbf{Tutorial Group:} 02\par
\vspace{4pt}

\begin{tabular}{@{}p{3.2cm}p{7cm}@{}}
Jaylen Avtarovski  & 24767939 \\[6pt]
Rhyse Williams & 24817532 \\[6pt]
Elina Kim & 24949885 \\[6pt]
\end{tabular}

\vspace{10pt}
\textbf{Declaration of Originality:}\par
\vspace{2pt}
{\small The work contained in this assignment, other than that specifically
attributed to another source, is that of the author(s). It is recognised
that, should this declaration be found to be false, disciplinary action
could be taken and the assignments of all students involved will be given
zero marks. In the statement below, I have indicated the extent to which I
have collaborated with other students, whom I have named.\par
\vspace{4pt}
Acknowledge use of any copied or reworked materials (eg.\ text, equations,
images, music, video) from generative artificial intelligence tools (GenAI)
(such as ChatGPT, DALL-E) by declaring its use (you'll also need to follow
the advice from your subject coordinator about whether you can use GenAI
tools in assessments for that subject). Find out more about how to reference
GenAI and ChatGPT in assessments.}\par

\vspace{6pt}
\textbf{GenAI Declaration:}\par
\vspace{2pt}
\todo{Declare any GenAI use for this report.}\par

\vspace{12pt}
\textbf{Signature} \,\,\,\,\,\,\,\,\, J.A. \,\,\,\,\,\,\,\,\,|\,\,\,\,\,\,\,\,\, E.K. \,\,\,\,\,\,\,\,\,|\,\,\,\,\,\,\,\,\, R.W. \hfill \textbf{Date} \todo{date}\par

\vspace{10pt}
\textbf{Statement of Collaboration:}\par
\vspace{2pt}
\todo{Describe how the team collaborated on Q1--Q3.}
\vspace{6pt}
\end{minipage}}

\vfill

%========================= CONTRIBUTION TABLE ==========================
\newpage

\begin{center}
\textbf{Demonstration Practice 2}\par
\vspace{4pt}
\begin{tabular}{|l|l|c|}
\hline
\textbf{Student Number} & \textbf{Student Name} &
  \textbf{\begin{tabular}{@{}c@{}}Overall Contribution Weighting\\ From 0 to 100\%\end{tabular}}\\
\hline
24767939 & Jaylen Avtarovski & 100\% \\ \hline
24817532 & Rhyse Williams    & 100\% \\ \hline
24949885 & Elina Kim         & 100\% \\ \hline
\end{tabular}
\end{center}

\vspace{6pt}
{\small Each member's individual mark is the group mark scaled by their
contribution level. If not indicated, all students get the same mark.}

\vspace{10pt}

\begin{center}\large\bfseries Answers to the Assignment Questions:\end{center}

% Marking checklist (applies to every question):
%  - comments using engineering concepts for every question (-2 per question)
%  - show all steps / MATLAB functions used for numerical results
%  - every figure numbered, captioned, referenced AND commented in the text (-2)
%  - axis labels with units, e.g. Time [s], \theta(t) [rad] (-2)
%  - several variables in one plot: different colours + legend (-2)
%  - zoomed plots for transients / steady-state error where relevant
%  - plots from MATLAB plot() (white background), no Simulink scope screenshots
%  - only your own figures (none from lecture notes/slides)

%=======================================================================
\section*{Q1: State-Space Model \hfill (40 marks)}
%=======================================================================

%-----------------------------------------------------------------------
\subsection*{Q1.1) System Modelling}
%-----------------------------------------------------------------------
% Derive from first principles (Newton / Kirchhoff / Lagrange).
% State n, m, p; give Ac, Bc, Cc, Dc. If nonlinear, Jacobian-linearise
% about a chosen equilibrium point (Remark 1).

Based on~\cite{aguilera2026}, the QUBE-Servo 2 is modelled as
\begin{subequations}\label{eq:ct_ss}
\begin{align}
  \dot{x}(t) &= A_c\,x(t) + B_c\,u(t),\\
  y(t)       &= C_c\,x(t) + D_c\,u(t),
\end{align}
\end{subequations}
with $n = \todo{n}$ states, $m = \todo{m}$ inputs and $p = \todo{p}$ outputs.

\todo{Derivation of the motor/load dynamics and the state vector $x(t)$.}

%-----------------------------------------------------------------------
\subsection*{Q1.2) System Parameters}
%-----------------------------------------------------------------------
% Use the datasheet values; if any are missing, identify them in MATLAB
% from experimental input-output data.

The physical parameters of the QUBE-Servo 2 are given in
Table~\ref{tab:params}, taken from the manufacturer's user
manual~\cite{qube_manual}.

\begin{table}[H]\centering
\caption{QUBE-Servo 2 system parameters~\cite{qube_manual}.}
\label{tab:params}
\begin{tabular}{llr}
\toprule
\textbf{Symbol} & \textbf{Description} & \textbf{Value}\\
\midrule
\multicolumn{3}{l}{\textit{DC motor}}\\
$R_m$ & Terminal resistance            & $8.4~\Omega$\\
$k_t$ & Torque constant                & $0.042~\mathrm{N\,m/A}$\\
$k_m$ & Motor back-emf constant        & $0.042~\mathrm{V/(rad/s)}$\\
$J_m$ & Rotor inertia                  & $4.0\times10^{-6}~\mathrm{kg\,m^2}$\\
$L_m$ & Rotor inductance               & $1.16~\mathrm{mH}$\\
$m_h$ & Module attachment hub mass     & $0.0106~\mathrm{kg}$\\
$r_h$ & Module attachment hub radius   & $0.0111~\mathrm{m}$\\
$J_h$ & Module attachment hub inertia  & $0.6\times10^{-6}~\mathrm{kg\,m^2}$\\
\multicolumn{3}{l}{\textit{Inertia disc module}}\\
$m_d$ & Disc mass                      & $0.053~\mathrm{kg}$\\
$r_d$ & Disc radius                    & $0.0248~\mathrm{m}$\\
\multicolumn{3}{l}{\textit{Rotary pendulum module}}\\
$m_r$ & Rotary arm mass                & $0.095~\mathrm{kg}$\\
$L_r$ & Rotary arm length              & $0.085~\mathrm{m}$\\
$m_p$ & Pendulum link mass             & $0.024~\mathrm{kg}$\\
$L_p$ & Pendulum link length           & $0.129~\mathrm{m}$\\
\multicolumn{3}{l}{\textit{Encoders and amplifier}}\\
      & Encoder resolution (quadrature) & $0.00307~\mathrm{rad/count}$\\
      & Recommended output voltage      & $\pm10~\mathrm{V}$\\
\bottomrule
\end{tabular}
\end{table}

\todo{Keep only the rows for the module used; add any identified parameters
(e.g.\ viscous friction) and the identification method.}

%-----------------------------------------------------------------------
\subsection*{Q1.3) System Model Discretisation}
%-----------------------------------------------------------------------
% Justify the sampling time, then c2d (ZOH).

A sampling time of $T_s = \todo{T_s}$ is selected because
\todo{justification, e.g.\ relative to the closed-loop bandwidth}. Using a
zero-order hold (\texttt{c2d}), the discrete-time model is
\begin{subequations}\label{eq:dt_ss}
\begin{align}
  x(k+1) &= A\,x(k) + B\,u(k),\\
  y(k)   &= C\,x(k) + D\,u(k).
\end{align}
\end{subequations}

\todo{Numerical $A$, $B$, $C$, $D$.}

%=======================================================================
\section*{Q2: State-Feedback Control (SFC) \hfill (40 marks)}
%=======================================================================

%-----------------------------------------------------------------------
\subsection*{Q2.1) Controllability}
%-----------------------------------------------------------------------
For this question, the controllability matrix of the continuous-time model
is
\begin{equation}
  \mathcal{C} = \begin{bmatrix} B_c & A_cB_c & \cdots & A_c^{n-1}B_c\end{bmatrix}.
\end{equation}
\todo{Numerical $\mathcal{C}$ (\texttt{ctrb}), its rank and justification.}

%-----------------------------------------------------------------------
\subsection*{Q2.2) State and Input References (Equilibrium Point)}
%-----------------------------------------------------------------------
Given the desired output reference $y^\star$, the steady-state references
$x^\star$ and $u^\star$ satisfy
\begin{equation}
  \begin{bmatrix} A - I & B\\ C & D\end{bmatrix}
  \begin{bmatrix} x^\star\\ u^\star\end{bmatrix}
  = \begin{bmatrix} 0\\ y^\star\end{bmatrix}.
\end{equation}
\todo{Solve for $x^\star$, $u^\star$ for your setup.}

%-----------------------------------------------------------------------
\subsection*{Q2.3) Closed-Loop Eigenvalues}
%-----------------------------------------------------------------------
% Project requirements: overshoot < 15 %, settling time < 5 s,
% steady-state accuracy < 50 (check units on Canvas).

The control requirements are an overshoot below $15\,\%$ and a settling time
below $5~\mathrm{s}$. The damping ratio and natural frequency follow from
\begin{equation}
  \zeta = \frac{-\ln(OS/100)}{\sqrt{\pi^2+\ln^2(OS/100)}},\qquad
  \omega_n = \frac{4}{\zeta\,t_s},
\end{equation}
and the discrete-time eigenvalues from $z_i = e^{s_i T_s}$.

\todo{Chosen $s_i$ and $z_i$ for all $n$ eigenvalues, with margin.}

%-----------------------------------------------------------------------
\subsection*{Q2.4) State-Feedback Gain}
%-----------------------------------------------------------------------
For the control law
\begin{equation}\label{eq:sfc}
  u(k) = -F\bigl(x(k) - x^\star\bigr) + u^\star,
\end{equation}
\todo{$F$ from \texttt{place(A,B,pz)} and verification of
$\operatorname{eig}(A-BF)$.}

%-----------------------------------------------------------------------
\subsection*{Q2.5) Closed-Loop Implementation}
%-----------------------------------------------------------------------
% run_observer = 0. Full-state feedback in a MATLAB Function block.
% Nonlinear systems: linear AND nonlinear model. Plot with
% plots_Linear_case.m / plots_Nonlinear_case.m, no scope screenshots.

\optfig{assets/images/fig_linear_sfc_response.png}{0.95}{Closed-loop response of the
linear model under state feedback.}{fig:q25_linear}

\todo{Comment on Figure~\ref{fig:q25_linear}: overshoot, settling time,
steady-state error and control effort versus requirements. Add zooms.}

%=======================================================================
\section*{Q3: Luenberger Observer Design \hfill (20 marks)}
%=======================================================================

The QUBE-Servo 2 measures only the two angles: the motor encoder gives
$\theta$ and the pendulum encoder gives $\alpha$, so
$y = C_c x = [\theta\;\;\alpha]^{\mathsf T}$. The velocities $\dot{\theta}$ and
$\dot{\alpha}$ used by the state-feedback law of Q2 are not measured. They could
be obtained by differencing the encoder signals, but at
$T_s = 1~\mathrm{ms}$ a single encoder count
($0.00307~\mathrm{rad} = 0.176^{\circ}$~\cite{qube_manual}) already appears as a
$176^{\circ}/\mathrm{s}$ velocity step. Instead, a Luenberger observer
reconstructs the full state from the measured outputs and the applied input,
using the model of Q1.

%-----------------------------------------------------------------------
\subsection*{Q3.1) Observability}
%-----------------------------------------------------------------------
With $n = 4$ states and $p = 2$ outputs, the observability matrix of the
continuous-time model~\eqref{eq:ct_ss} is the $np \times n = 8\times 4$ matrix
(MATLAB: \texttt{OM = obsv(Ac,Cc)})
\begin{equation}\label{eq:obsv}
  \mathcal{O} =
  \begin{bmatrix} C_c\\ C_cA_c\\ C_cA_c^{2}\\ C_cA_c^{3}\end{bmatrix}
  =
  \left[\begin{array}{rrrr}
     1 & 0 & 0 & 0\\
     0 & 1 & 0 & 0\\ \hline
     0 & 0 & 1 & 0\\
     0 & 0 & 0 & 1\\ \hline
     0 & 54.236 & -6.179 & -1.786\\
     0 & 167.67 & -6.107 & -5.521\\ \hline
     0 & -634.55 & 49.086 & 75.128\\
     0 & -1256.9 & 71.453 & 209.06
  \end{array}\right].
\end{equation}

The first four rows of~\eqref{eq:obsv} form the $4\times4$ identity matrix,
because $C_c$ selects the two angles and $C_cA_c$ selects the two angular
velocities. Hence
\begin{equation}
  \operatorname{rank}(\mathcal{O}) = 4 = n,
\end{equation}
confirmed with \texttt{rank(OM)}, and the system is \textbf{observable}.

From a control viewpoint, every state shows up in the measured outputs or in
their time derivatives. The velocities are exactly the rates of change of the
measured angles, so the initial state can be reconstructed from a finite window
of input-output data. Rank~4 is reached with only the first time derivative,
which means the observer only has to infer the velocities. The ZOH
discretisation does not change this: \texttt{rank(obsv(A,C))} $= 4$ as well,
because $T_s$ is far below the time scale of every open-loop mode (the open-loop
eigenvalues are $0$, $9.95$, $-4.02$ and $-17.63~\mathrm{rad/s}$), so
pathological sampling cannot occur.

%-----------------------------------------------------------------------
\subsection*{Q3.2) Observer Eigenvalues}
%-----------------------------------------------------------------------
By the separation principle~\cite{aguilera2026}, the eigenvalues of the
observer-based closed loop are the union of the controller eigenvalues,
$\operatorname{eig}(A-BF)$, and the observer eigenvalues,
$\operatorname{eig}(A-LC)$. The observer can therefore be designed
independently. It should be clearly faster than the controller, so that the
estimation error has vanished before it can noticeably affect the transient set
by $F$.

The state-feedback eigenvalues from Q2.3 are
$s_{1,2} = -2 \pm j2.097$ (dominant pair: $\zeta = 0.69$,
$\omega_n = 2.90~\mathrm{rad/s}$) and $s_3 = -30$, $s_4 = -35~\mathrm{rad/s}$.
The observer eigenvalues are placed about three times faster than the fastest
controller eigenvalue:
\begin{equation}\label{eq:obs_poles_s}
  s_{o} = \{-100,\; -110,\; -120,\; -130\}~\mathrm{rad/s}.
\end{equation}
The reasons for this choice are:
\begin{itemize}[nosep]
  \item The slowest observer mode has a time constant
        $\tau = 1/100 = 10~\mathrm{ms}$, so the estimation error settles
        ($\approx 4\tau$) in about $40~\mathrm{ms}$. That is $50$ times faster
        than the dominant controller pair and about $54$ times shorter than the
        $2.15~\mathrm{s}$ closed-loop settling time.
  \item Real eigenvalues give a non-oscillatory error decay. They are distinct
        because \texttt{place} only allows each eigenvalue a multiplicity of
        up to $\operatorname{rank}(C) = 2$.
  \item The observer is not made faster still. Faster eigenvalues need a larger
        gain $L$, which passes more encoder quantisation noise into $\hat{x}$
        and increases the initial ``peaking'' of the velocity estimates
        (Figure~\ref{fig:q34_estimation}). At $|s_o|T_s \le 0.13$, each
        observer mode still spans more than seven samples, so the discrete
        design is well resolved.
\end{itemize}

Mapping these to the $z$-plane with $z = e^{sT_s}$, $T_s = 1~\mathrm{ms}$
(MATLAB: \texttt{po\_z = exp(po\_s*ts)}), gives the discrete-time observer
eigenvalues
\begin{equation}\label{eq:obs_poles_z}
  z_{o} = \{0.9048,\; 0.8958,\; 0.8869,\; 0.8781\}.
\end{equation}
All lie on the positive real axis, inside the unit circle and closer to the
origin than the controller eigenvalues
$\{0.9980 \pm j0.0021,\; 0.9704,\; 0.9656\}$. This confirms in the $z$-domain
that the estimation error decays faster than the regulated response.

%-----------------------------------------------------------------------
\subsection*{Q3.3) Observer Gain}
%-----------------------------------------------------------------------
The discrete-time Luenberger observer is
\begin{equation}\label{eq:observer}
  \hat{x}(k+1) = A\hat{x}(k) + Bu(k) + L\bigl(y(k)-\hat{y}(k)\bigr),
  \qquad \hat{y}(k) = C\hat{x}(k),
\end{equation}
i.e.\ a copy of the plant model driven by the same input and corrected by the
output innovation $y(k)-\hat{y}(k)$. Subtracting~\eqref{eq:observer} from the
plant model~\eqref{eq:dt_ss}, the estimation error
$e(k) = x(k) - \hat{x}(k)$ evolves as
\begin{equation}\label{eq:err_dyn}
  e(k+1) = (A - LC)\,e(k),
\end{equation}
independent of $u(k)$. Because
$\operatorname{eig}(A-LC) = \operatorname{eig}(A^{\mathsf T}-C^{\mathsf T}L^{\mathsf T})$,
choosing $L$ is the dual of a state-feedback design for the pair
$(A^{\mathsf T}, C^{\mathsf T})$. It is solvable because $(A, C)$ is observable
(Q3.1). In MATLAB:
\begin{lstlisting}
po_z = exp([-100 -110 -120 -130]*ts);
Lo   = place(A', C', po_z)';
eig(A - Lo*C)
\end{lstlisting}
which gives the $n\times p = 4\times2$ observer gain
\begin{equation}\label{eq:L}
  L =
  \begin{bmatrix}
     0.2108  & -0.0104\\
    -0.0152  &  0.2120\\
    10.466   & -1.188\\
    -2.168   & 10.861
  \end{bmatrix}.
\end{equation}
Evaluating the eigenvalues of the error-dynamics matrix gives
\begin{equation}
  \operatorname{eig}(A - LC) = \{0.9048,\; 0.8958,\; 0.8869,\; 0.8781\},
\end{equation}
identical to~\eqref{eq:obs_poles_z}, so the observer has the desired
eigenvalues.

The structure of $L$ has a clear physical reading. Column~1 weights the
$\theta$ innovation $\theta - \hat{\theta}$, and column~2 weights the $\alpha$
innovation $\alpha - \hat{\alpha}$. Each angle estimate is corrected mainly by
its own measurement ($0.21$ per sample, on the diagonal of the first two
rows). Each velocity estimate is corrected by about $10.5~\mathrm{s^{-1}}$ per
radian of its own angle innovation, which works like a model-based, filtered
differentiator. The small off-diagonal terms account for the coupling between
the arm and the pendulum in $A$.

%-----------------------------------------------------------------------
\subsection*{Q3.4) Closed-Loop Implementation}
%-----------------------------------------------------------------------
The observer was implemented in the provided Simulink models
(\path{sim_DP2_2026_SFC_Linear.slx} and
\path{sim_DP2_2026_SFC_Nonlinear.slx}) inside the MATLAB Function block
\emph{Digital Controller}, which is sampled at $T_s = 1~\mathrm{ms}$. The
controller now feeds back the estimated state:
\begin{equation}\label{eq:sfc_obs}
  u(k) = -F\bigl(\hat{x}(k) - x^\star\bigr) + u^\star.
\end{equation}
The core of the block (with \texttt{run\_observer = 1}) is:
\begin{lstlisting}
x_hat = x_hat_k;                            % estimate used at sample k
xu_ref = [A-eye(n) B; C zeros(p,m)] \ [zeros(n,1); y_ref];
x_ref = xu_ref(1:n);  u_ref = xu_ref(n+1:n+m);   % Q2.2
u_k = -F*(x_hat_k - x_ref) + u_ref;       
u_k = min(max(u_k, -10), 10);               % +/-10 V amplifier limit
y_hat   = C*x_hat_k;                        % Luenberger observer (Q3.3)
x_hat_k = A*x_hat_k + B*u_k + Lo*(y_k - y_hat);
\end{lstlisting}
The observer is updated with the \emph{saturated} voltage that is actually
applied to the motor. Feeding it the unsaturated command would make its model
input differ from the plant input whenever the amplifier limits, and the
estimate would drift.

The controller was applied to both the linearised model (Q1.1) and the
nonlinear rotary-pendulum model. The nonlinear model is the same
Euler--Lagrange model before linearisation, with
$\sin\alpha$, $\cos\alpha$, $\dot{\theta}^2$, $\dot{\alpha}^2$ and
$\dot{\theta}\dot{\alpha}$ terms retained, implemented in
\path{rip_nonlinear_dynamics.m}. Both runs use the arm reference
$\theta^\star = 30^{\circ}$ applied as a step at $t = 1~\mathrm{s}$ and
$\alpha^\star = 0$ (upright), for which Q2.2 gives
$x^\star = [\theta^\star\;0\;0\;0]^{\mathsf T}$ and $u^\star = 0$. As Remark~5
asks, the plant starts at rest upright, $x(0) = 0$, while the observer is
deliberately initialised with
\begin{equation}
  \hat{x}(0) = \bigl[\,10^{\circ}\;\; -5^{\circ}\;\; 0\;\; 0\,\bigr]^{\mathsf T}
  \neq x(0).
\end{equation}
All signals were exported to the workspace and plotted with
\path{plot_closed_loop.m}. The script \path{run_all_cases.m}
regenerates every figure.

\textbf{Observer convergence.}
Figure~\ref{fig:q34_estimation} compares the true and estimated states of the
nonlinear model, with a zoom on the first $0.1~\mathrm{s}$.
Figure~\ref{fig:q34_error} shows the estimation error on a logarithmic scale.
The angle estimates converge onto the measured angles within a few
milliseconds, and both angle errors stay below $0.1^{\circ}$ after
$54~\mathrm{ms}$ in both the linear and nonlinear cases. This agrees with the
$\approx 40~\mathrm{ms}$ predicted from the slowest observer eigenvalue in Q3.2.
After the first few samples, the four error curves in
Figure~\ref{fig:q34_error} are straight lines on the logarithmic axis, which
is the exponential decay of~\eqref{eq:err_dyn}. The final error is about
$10^{-12}$, at the level of numerical round-off.

The zoom in Figure~\ref{fig:q34_estimation} also shows the price of a fast
observer. The velocities are not measured, so the initial $10^{\circ}$ and
$-5^{\circ}$ angle innovations are converted through the large velocity gains of
$L$ into velocity-estimate errors that peak at $389^{\circ}/\mathrm{s}$
($\dot{\theta}$) and $254^{\circ}/\mathrm{s}$ ($\dot{\alpha}$) about
$8$--$9~\mathrm{ms}$ after start-up. This is the typical peaking of a
high-gain observer, and it is why the observer eigenvalues were not pushed
further left.

\textbf{Closed-loop performance.}
Figures~\ref{fig:q34_closed_loop_L} and~\ref{fig:q34_closed_loop_NL} show the
response of the observer-based controller on the linear and nonlinear models.
Table~\ref{tab:q34_metrics} compares it with the full-state feedback of Q2.5.
Because $u(k)$ is computed from $\hat{x}(k)$, the wrong initial estimate acts
like an initial-condition disturbance. At $k = 0$ the controller ``sees'' the
pendulum $5^{\circ}$ off upright and commands $u(0) = 8.98~\mathrm{V}$. The
voltage then saturates at $-10~\mathrm{V}$ for $21~\mathrm{ms}$ while the
estimate converges. This briefly kicks the pendulum to $|\alpha| = 5.8^{\circ}$,
compared with $1.5^{\circ}$ under full-state feedback, and moves the arm by
about $12^{\circ}$ before the controller recovers. Under full-state feedback,
nothing happens before the step.

Once the estimation error has converged (well before the step at
$t = 1~\mathrm{s}$), $\hat{x}(k) \approx x(k)$ and, by the separation
principle, the response is governed only by $\operatorname{eig}(A - BF)$. The
step response is accordingly almost identical to Q2.5: $4.6\,\%$ overshoot
(requirement $< 15\,\%$), $2.16~\mathrm{s}$ settling time to the $2\,\%$ band
(requirement $< 5~\mathrm{s}$) and negligible steady-state error. The slightly
lower overshoot than in Q2.5 comes from the arm still returning from the
start-up transient when the step is applied. All requirements are met in both
models. Saturation briefly breaks the linearity that the separation principle
assumes. Because the observer uses the applied (saturated) input, the
estimation-error dynamics~\eqref{eq:err_dyn} are unaffected and the estimate
stays correct throughout.

\begin{table}[H]\centering
\caption{Simulated step-response metrics ($\theta^\star = 30^{\circ}$ at
$t = 1~\mathrm{s}$; overshoot and $2\,\%$ settling time measured from the
step).}
\label{tab:q34_metrics}
\begin{tabular}{llcccc}
\toprule
\textbf{Model} & \textbf{Feedback} & \textbf{OS [\%]} & $\bm{t_s}$ \textbf{[s]} &
  $\bm{\max|\alpha|}$ \textbf{[deg]} & $\bm{\max|V_m|}$ \textbf{[V]}\\
\midrule
Linear    & $x(k)$ (Q2.5)       & 5.32 & 2.15 & 1.47 & 2.18\\
Nonlinear & $x(k)$ (Q2.5)       & 5.30 & 2.15 & 1.47 & 2.18\\
Linear    & $\hat{x}(k)$ (Q3.4) & 4.64 & 2.16 & 5.82 & 10.0 (sat.)\\
Nonlinear & $\hat{x}(k)$ (Q3.4) & 4.59 & 2.16 & 5.84 & 10.0 (sat.)\\
\midrule
\multicolumn{2}{l}{Requirement} & $< 15$ & $< 5$ & -- & $\le 10$\\
\bottomrule
\end{tabular}
\end{table}

The linear and nonlinear responses differ by less than $0.1$ percentage points
of overshoot. The pendulum never leaves $|\alpha| < 6^{\circ}$, where
$\sin\alpha \approx \alpha$ is accurate to better than $0.2\,\%$, so the
linearised model on which both $F$ and $L$ are based remains a faithful
description of the plant.

On the real QUBE-Servo 2, two practical points follow from these results.
First, the start-up kick can be avoided by initialising $\hat{\theta}(0)$ and
$\hat{\alpha}(0)$ with the first encoder readings, since only the
unmeasured velocities then need to be estimated. Second, the simulated
steady-state error is zero only because the model is exact. Unmodelled
Coulomb friction and parameter mismatch on the hardware would leave an offset
in $\theta$ that this proportional state feedback cannot remove without
integral action.

\optfig{assets/images/fig_nonlinear_obs_estimation.png}{1.0}{True ($x$, solid)
and estimated ($\hat{x}$, dashed) states of the nonlinear model with
$\hat{x}(0) = [10^{\circ}\; -5^{\circ}\; 0\; 0]^{\mathsf T} \neq x(0) = 0$.
Left: full simulation; right: zoom on the first $0.1~\mathrm{s}$ showing
observer convergence.}{fig:q34_estimation}

\optfig{assets/images/fig_linear_obs_error.png}{0.9}{Magnitude of the
estimation error $e = x - \hat{x}$ (linear model, logarithmic scale) during the
first $0.2~\mathrm{s}$; the straight-line decay reflects the exponential error
dynamics $e(k+1) = (A - LC)e(k)$.}{fig:q34_error}

\optfig{assets/images/fig_linear_obs_response.png}{0.85}{Observer-based state
feedback on the linear model: arm angle $\theta$ with reference $\theta^\star$
(dotted lines: $\pm2\,\%$ band), pendulum angle $\alpha$, and motor voltage
$V_m$.}{fig:q34_closed_loop_L}

\optfig{assets/images/fig_nonlinear_obs_response.png}{0.85}{Observer-based state
feedback on the nonlinear model: arm angle $\theta$ with reference $\theta^\star$
(dotted lines: $\pm2\,\%$ band), pendulum angle $\alpha$, and motor voltage
$V_m$.}{fig:q34_closed_loop_NL}

\clearpage

%=======================================================================
\begin{thebibliography}{9}

\bibitem{aguilera2026}
R. P. Aguilera, \emph{Lecture Notes on State-Feedback Control}, University of
Technology Sydney, Australia, 2026.

\bibitem{ctms}
B. Messner and D. Tilbury, \emph{Control Tutorials for MATLAB and Simulink},
University of Michigan, USA. [Online]. Available:
\url{http://ctms.engin.umich.edu/CTMS}

\bibitem{qube_manual}
Quanser Inc., \emph{QUBE-Servo 2 User Manual}, v1.0, Markham, ON, Canada.

\end{thebibliography}

\end{document}
